\documentclass{article}
\usepackage[T1]{fontenc}
\usepackage{lmodern}
\usepackage{graphicx}
\usepackage{amsmath,amssymb}
\usepackage[a4paper,margin=2cm]{geometry}
\usepackage[absolute,overlay]{textpos}
\setlength{\TPHorizModule}{1mm}
\setlength{\TPVertModule}{1mm}
\usepackage[indonesian]{babel}
\usepackage{hyperref}
\setlength{\emergencystretch}{2em}
% unit: Halaman 13

\begin{document}

% === Halaman 13 (python/native) ===

Rinaldi Munir/II4021 Kriptografi

13

Contoh: F23 mempunyai anggota \{0, 1, 2, …, 22\}. 

      Contoh operasi aritmetika di dalam F23:


12 + 20 = 9 (karena 12 + 20 = 32 mod 23 = 9)


8  9 = 3       (karena 8  9 = 72 mod 23 = 3)

\end{document}
